%!TEX root = ../main.tex
% =============================================================
%  CAPÍTULO 1 — INTRODUCCIÓN            (Entrega 1 · 5 % · 9/6)
% =============================================================

\chapter{Introducción}\label{cap:introduccion}

\begin{guia}[Guía de la cátedra: qué se evalúa en este capítulo]
\begin{enumerate}
    \item \textbf{Descripción del problema}: concreto, anclado en evidencia
          (datos duros con fuente, casos reales, competidores nombrados).
    \item \textbf{Objetivo general}: uno solo, medible, en un párrafo.
          Los \textbf{objetivos específicos} se desprenden de él y
          \emph{no} son los pasos del plan de trabajo.
    \item \textbf{Hipótesis}: explícita, formulada como pregunta o
          afirmación de investigación que pueda responderse en el
          capítulo~\ref{cap:conclusiones}.
    \item \textbf{Relevancia}: por qué importa, a quién beneficia, qué les
          falta a las soluciones existentes.
    \item \textbf{Alcance y recursos}: qué entra y qué no; stack y fuentes
          de datos previstos. El ``alcance que nunca cierra'' es el riesgo
          número uno de la materia.
\end{enumerate}
Extensión orientativa: 4 a 6 páginas. Redacción impersonal, sin ``yo
creo'', oraciones de 15 a 22 palabras, sin palabras vacías
(``claramente'', ``obviamente''). Toda cifra con su fuente citada, por
ejemplo \citep{breiman2001}.
\end{guia}

\section{Problema y contexto}

En Argentina, los motociclistas concentran una porción desproporcionada
de las víctimas fatales de tránsito. Datos oficiales de 2025 de la
Agencia Nacional de Seguridad Vial (ANSV) ubicaron a los motociclistas
en el 46\% de las víctimas fatales de siniestros viales del país
\citep{infoecos2026}, cifra consistente con la reportada
por coberturas periodísticas independientes para el mismo período, en un
rango de 46-47\% \citep{lanacion2025,datoposta2025}. El informe Global
status report on road safety 2023 de la Organización Mundial de la Salud
confirma el patrón a escala global: los motociclistas, junto con
peatones y ciclistas, concentran una fracción desproporcionada de las
muertes de tránsito, sobre todo en países de ingresos medios como
Argentina \citep{who2023}.

Ese riesgo conocido convive con una segunda dificultad, menos discutida:
la información pública que podría orientar una decisión de circulación
está fragmentada. Los registros históricos de siniestros, los reportes
de inseguridad, el clima, el estado de una vía y las noticias locales
recientes se producen y publican de manera independiente, cada uno con
su propia cobertura geográfica, su propia vigencia y su propio nivel de
verificación. Un motociclista que quiere saber qué antecedentes existen
en una zona antes de iniciar un recorrido no tiene, hoy, un punto único
que se lo muestre sin mezclar esas procedencias entre sí.

MotoSafe GBA no se conformará con dejar expuesta esa fragmentación: la
reducirá de forma activa incorporando reportes colaborativos de
motociclistas y noticias verificadas a medida que se sumen usuarios, y
hará explícito qué zonas siguen con cobertura baja mientras esa
incorporación avanza. El sistema buscará acercarse a una estimación útil
de la gravedad esperada de un siniestro, sin convertirla en una promesa
de seguridad para un trayecto puntual: cada capa mostrará su
procedencia, antigüedad y cobertura junto con la información misma, en
lugar de resumir todo en una única cifra que oculte esas diferencias.

\section{Objetivo general y objetivos específicos}

Objetivo general: diseñar e implementar un MVP web que permita
consultar, visualizar y comparar información de riesgo vial y contexto
reciente para motociclistas del AMBA, con trazabilidad de fuentes,
cobertura declarada y límites de interpretación explícitos.

Objetivos específicos:

\begin{itemize}
\item
  Integrar datos históricos de siniestralidad de motociclistas en CABA y
  representarlos como zonas de referencia sobre un mapa interactivo.
\item
  Implementar un modelo de clasificación de gravedad condicionado a la
  ocurrencia de un siniestro, documentando de forma explícita qué estima
  y qué no estima.
\item
  Incorporar capas de inseguridad general, noticias de prensa, clima y
  reportes colaborativos sin equiparar su procedencia ni su cobertura
  con la de los datos históricos de siniestros.
\item
  Diseñar una interfaz cartográfica accesible, operable en escritorio y
  en dispositivos móviles.
\item
  Construir una arquitectura reproducible, con API propia, base
  PostgreSQL/PostGIS, proceso de ingesta de prensa programado y una vía
  de despliegue en la nube documentada.
\end{itemize}

\section{Relevancia y beneficiarios}

El aporte de MotoSafe GBA estará en organizar información pública
dispersa y reciente en una interfaz orientada a un usuario concreto,
haciendo visibles la procedencia, la antigüedad y la cobertura de cada
fuente antes de que esa información se use para interpretar una zona o
comparar una ruta.

Los beneficiarios directos serán los motociclistas del AMBA que hoy
deben cruzar manualmente registros oficiales, prensa y reportes de
terceros para formarse una idea de una zona. De forma indirecta, la
arquitectura documentada del proyecto (contrato del modelo, ficha de
datos, protocolo de cobertura) también servirá como antecedente
metodológico para trabajos posteriores que integren fuentes heterogéneas
de riesgo vial sin diluir sus diferencias de procedencia.

\section{Delimitación, alcance y barreras}

El trabajo se concentrará en un MVP web. La evidencia histórica que
sostenga el modelo de gravedad se validará sobre la Ciudad Autónoma de
Buenos Aires; el resto del Gran Buenos Aires se incorporará como capa
exploratoria, sin la misma validación. El proyecto no incluirá
aplicación móvil nativa, gestión de flotas ni garantía de seguridad para
un trayecto. El alcance incluye autenticación de identidad de usuarios,
con registro, verificación por código y recuperación de contraseña.

El proyecto enfrentará barreras de distinto tipo. La naturaleza temporal
del fenómeno exigirá validar el modelo contra datos de un período no
visto durante el entrenamiento, un protocolo más exigente que una
partición aleatoria y sujeto a la disponibilidad de datos históricos
completos. Integrar prensa, reportes colaborativos y datos oficiales sin
equiparar su procedencia ni su cobertura tampoco se resolverá de una
vez: va a exigir un trabajo de curación con medición continua. El
seguimiento de estas barreras se centralizará en el registro de gestión
del proyecto.

\section{Pregunta de investigación e hipótesis de trabajo}

Pregunta de investigación: ¿puede una plataforma que integre
siniestralidad histórica, inseguridad, prensa y clima en una interfaz
cartográfica mejorar la información disponible para un motociclista sin
sobreestimar su capacidad predictiva?

Hipótesis de trabajo: una interfaz que declare qué sabe, desde cuándo lo
sabe y con qué cobertura lo sabe es más útil y más defendible, para un
motociclista, que un único indicador de riesgo que mezcle fuentes
heterogéneas sin distinguirlas. La utilidad de MotoSafe GBA está en esa
trazabilidad, no en predecir con precisión la ocurrencia de un siniestro
individual.

\section{Enfoque y recursos}

El enfoque avanzará por etapas de relevamiento, diseño, implementación y
evaluación, priorizando en cada una que la procedencia y la cobertura de
los datos queden explícitas antes de integrarse a una consulta única.

Los recursos del proyecto se agruparán en cuatro tipos. Humanos: un
único desarrollador a cargo de todo el proyecto. Tecnológicos: el
repositorio de control de versiones Git como fuente única de verdad, un
entorno de contenedores Docker para reproducibilidad, y dos asistentes
de inteligencia artificial usados como apoyo, uno para auditoría de
código, implementación técnica y documentación operativa, y otro para
redacción académica, auditoría cruzada de la documentación y
verificación de coherencia con los requisitos de la cátedra. De
información: los datasets públicos de siniestralidad e inseguridad, y
las fuentes de prensa que alimentan el sistema, sin costo de
adquisición. Económicos: el desarrollo se sostendrá con recursos propios
y servicios gratuitos o de bajo costo.

El cronograma de entregas seguirá el esquema de diez entregas
establecido por la cátedra para el Proyecto Final de Ingeniería en
Informática, desde la idea y el abstract iniciales hasta la presentación
final con el proyecto funcionando, pasando por los capítulos, el MVP y
el prototipo intermedios.
